% BAN TIENG VIET cua 05_design_results.tex (2026-10-08). So lieu giu nguyen.
\section{Thiết kế thực nghiệm}
\label{sec:design}

Thiết kế này so sánh bốn tầng phòng thủ với chín loại nhật ký trên một thửa
ruộng cố ý vi phạm mọi giả định của mô hình sinh đôi số: khác với đánh giá
trước đây, mà $0/60$ ca kết tội oan chỉ đo tính tự nhất quán (Phụ lục~S10.4),
ruộng ở đây được sinh ra bởi một tiến trình riêng biệt, khác mô hình sinh đôi
số theo bảy trục, mỗi trục đều có nguồn đã công bố (mọi giá trị ở Bảng~S6 của
Phụ lục; danh sách từng trục và ngữ nghĩa ghi chép theo sự kiện ở S20.1). Thời
tiết là ERA5 theo giờ (kho lưu trữ Open-Meteo) cho Cần Thơ, Sóc Trăng và Cà Mau
trong năm 2024 (khan nước) và 2025 (thừa nước), cửa sổ từ ngày 1 tháng Ba đến
ngày 29 tháng Năm ($2\,160$~giờ), mùa khô mà AWD được áp dụng. Tầng~1 là bộ
phát hiện mô hình điểm đã công bố (hồ sơ mặc định, ngưỡng tuyệt đối
$\tau_0=5$~mm); Tầng~2 là bộ phát hiện bền vững của \eqref{eq:robust} với
ngưỡng tỉ lệ của \eqref{eq:tau}, chạy trên cả $\Theta$ rút gọn ($24$ lần phát
lại mỗi nhật ký, $10$ seed) và $\Theta$ đầy đủ ($150$ lần phát lại, dùng kiểm
chéo); Tầng~3 là neo thiết bị của \eqref{eq:anchor}; Tầng~4 chỉ có vật lý, dùng
làm mốc. Trong chín loại nhật ký, ba loại là giả mạo ngây thơ và bốn loại là
đối kháng, được dựng bằng một bộ tối ưu chứ không theo một công thức cố định
(danh sách đầy đủ ở S16); mọi thống kê tấn công đều có điều kiện trên cờ
\texttt{changed}, vì một đối thủ không tìm được nước đi cải thiện sẽ trả về
nguyên bản nhật ký, và các phép so sánh dùng mốc của chính biến thể đó (hai lỗi
của bộ khung thử nghiệm, S10.2). Với mỗi tầng chúng tôi báo cáo tỷ lệ lọt lưới
(nhật ký giả mà không tầng vật lý nào đánh dấu) và tỷ lệ kết tội oan trên nhật
ký trung thực; về tác động tín chỉ, báo cáo $n_{\mathrm{dry}}$ của
\eqref{eq:ndry} so với mốc ruộng trung thực, theo ghép cặp. Tổng lượng nước chỉ
xuất hiện như một chỉ số chẩn đoán (S10.3).

\section{Kết quả}
\label{sec:results}

\subsection{Bộ phát hiện đã công bố không an toàn, chứ không chỉ lạc quan}

Bảng~\ref{tab:falseacc} đưa ra kết quả phủ định trung tâm: bộ phát hiện mô hình
điểm không kết tội ai trên nhật ký điều kiện giấy, tái hiện đúng $0/60$ của lần
trước, nhưng kết tội oan $26{,}7\%$ nông dân trung thực trên nhật ký điều kiện
ruộng. Nguyên nhân là \eqref{eq:gap} với một ngưỡng tuyệt đối --- quên ghi, làm
tròn hoặc ghi trễ một ngày đều tạo ra $G>5$~mm mà không có hành vi sai trái
nào; chênh lệch đo được trên nhật ký trung thực điều kiện ruộng lên tới
$65$~mm. Ngưỡng tỉ lệ của \eqref{eq:tau} hạ con số này xuống $5{,}0\%$, và neo
thiết bị hạ xuống $0/60$.

\subsection{Vật lý không nhìn thấy tấn công sinh lợi}

Bảng~\ref{tab:credit} là kết quả phủ định thứ hai, và là kết quả quan trọng về
mặt kinh tế. Mọi con số thổi phồng đều là so sánh \emph{ghép cặp} với nhật ký
trung thực của cùng quỹ đạo, chỉ tính trên các tấn công thật sự thay đổi bản
ghi; cách tính không ghép cặp sẽ phóng đại lợi ích lên khoảng hai lần (S15).
Dịch thời điểm nâng số pha khô được cấp tín chỉ từ $4{,}78$ lên $6{,}33$ mỗi
vụ, tức $+33\%$ trên chín nhật ký bị sửa ($t=5{,}29$); với tập bất định đầy đủ,
$+38\%$ trên ba nhật ký ($t=3{,}46$); đối thủ duy lý đạt $+35\%$ và $+40\%$.
Gộp khai báo hầu như không được gì khi ghép cặp ($+1\%$, $+2\%$), và hai kiểu
giả mạo thô sơ thậm chí \emph{làm giảm} số pha khô được cấp tín chỉ ($-1\%$,
$-17\%$): chỉ có phép hoán vị dấu thời điểm giữ nguyên tổng lượng nước là sinh
lợi, đúng như Mệnh đề~\ref{prop:coincide} dự đoán. Tỷ lệ lọt lưới trong
Bảng~\ref{tab:defence} là phần cốt lõi của kết quả: dịch thời điểm lọt qua tầng
vật lý bền vững ở $8$ trong $9$ nhật ký bị sửa, đối thủ duy lý ở $6$ trong $6$.
Mô hình điểm bắt được nhiều hơn ($4$ trong $9$), nhưng nó lại chính là bộ phát
hiện kết tội oan $26{,}7\%$ nông dân trung thực; tầng bền vững hạ tỷ lệ đó xuống
$5{,}0\%$ và đổi lại giao phần lớn khả năng phát hiện về tay đối thủ ($1$ trong
$9$) --- \emph{không cấu hình vật lý nào thoát khỏi sự đánh đổi đó}, vì cả hai
tầng đều là hàm của quỹ đạo phát lại. Bộ lọc thấm $V_5$ kích hoạt $648$ lần,
nhưng chỉ trên kiểu giả mạo thổi phồng độ sâu thô sơ: hữu ích, và trực giao với
mối đe doạ sinh lợi.

\subsection{Neo thiết bị bịt mọi lớp, và cái giá của sự bền vững}

Bảng~\ref{tab:defence} cũng đưa ra nửa kiến tạo: trong $281$ nhật ký được kiểm
toán, $116$ nhật ký thật sự bị sửa, và neo thiết bị bắt được $116/116$ trong
khi không kết tội oan nhật ký nào trong $60$ nhật ký trung thực điều kiện ruộng
($37/37$ và $0/18$ với $\Theta$ đầy đủ). Giả mạo ngây thơ bị mọi tầng bắt được;
khác biệt nằm hoàn toàn ở các dòng đối kháng, nơi cả vật lý lẫn tổng lượng đều
thất bại và chỉ có neo thiết bị đứng vững. Tính bền vững không miễn phí: trên
chín nhật ký gian lận bỏ sót thật sự bị sửa, mô hình điểm bắt được $9/9$ trong
khi quy tắc bền vững bắt $4/9$ ($4/4$ so với $2/4$ ở $\Theta$ đầy đủ) --- một sự
mất mát năng lực phát hiện, không phải một cải thiện Pareto, dù chúng tôi coi
đó là phía đúng của sự đánh đổi. Mọi kết luận đều đúng với cả $\Theta$ rút gọn
($24$ lần phát lại, $281$ nhật ký, $10$ seed) và $\Theta$ đầy đủ ($150$ lần phát
lại, $85$ nhật ký, $3$ seed): kết tội oan $0/60$ và $0/18$, tính đầy đủ của neo
$116/116$ và $37/37$, lọt lưới dịch thời điểm $8/9$ và $3/3$, thổi phồng ghép
cặp $+33\%$ và $+38\%$ (Phụ lục~S17).
